\slajdid{10-s039}
\begin{frame}[label=10-s039]{Diodna logika}
\gusttekst\setlength{\abovedisplayskip}{3pt}\setlength{\belowdisplayskip}{3pt}\setlength{\abovedisplayshortskip}{2pt}\setlength{\belowdisplayshortskip}{2pt}Diode su „dobri“ elementi elektronike koji mogu na lak način da obezbede višulaznu logiku.\par\bigskip Diodna ILI funkcija
\begin{columns}[c,onlytextwidth]\column{.29\textwidth}\centering\begin{tikzpicture}[x=.62cm,y=.8cm,font=\sitneoznake,>=Latex]\ctikzset{bipoles/length=.6cm}\draw(-1.8,1)node[ocirc]{}node[left]{$V_{I1}$}--(-1,1);\begin{scope}[shift={(0,1)}]\draw(-1,0)--(-.3,0);\draw(-.3,-.3)--(-.3,.3)--(.3,0)--cycle;\draw(.3,0)--(1,0);\draw(.3,-.3)--(.3,.3);\end{scope}\draw(1,1)--(1.4,1);\node[below]at(0,0.85){D1};\draw[<->](-.5,1.55)--(.5,1.55);\draw(-.5,1.4)--(-.5,1.7);\draw(.5,1.4)--(.5,1.7);\node[above]at(0,1.7){$+\ V_{D1}$};\draw[->](-1.7,1.35)--(-1.1,1.35)node[midway,above]{$I_{D1}$};\draw(-1.8,-1)node[ocirc]{}node[left]{$V_{I2}$}--(-1,-1);\begin{scope}[shift={(0,-1)}]\draw(-1,0)--(-.3,0);\draw(-.3,-.3)--(-.3,.3)--(.3,0)--cycle;\draw(.3,0)--(1,0);\draw(.3,-.3)--(.3,.3);\end{scope}\draw(1,-1)--(1.4,-1);\node[below]at(0,-1.15){D2};\draw[<->](-.5,-0.44999999999999996)--(.5,-0.44999999999999996);\draw(-.5,-0.6)--(-.5,-0.30000000000000004);\draw(.5,-0.6)--(.5,-0.30000000000000004);\node[above]at(0,-0.30000000000000004){$+\ V_{D2}$};\draw[->](-1.7,-0.65)--(-1.1,-0.65)node[midway,above]{$I_{D2}$};\draw(1.4,1)--(1.4,-1);\draw(1.4,0)--(2.8,0)node[ocirc]{}node[right]{$V_O$};\draw(2.1,0)to[R,l=$R$](2.1,-1.8)node[ground]{};\draw[->](1.6,-.5)--(1.6,-1.1);\node[left]at(1.6,-1.55){$I_R$};\end{tikzpicture}
\column{.69\textwidth}
Diode su katodama povezane u zajedničku tačku i preko puldaun otpornika povezane na masu. Posmatrajmo situaciju da su oba ulazna napona $V_{I1}=V_{I2}=0$. Ako bi pretpostavili da bilo koja od dioda vodi napon na njoj bi bio $V_D$ i po konturi na primer diode D1 i otpornika R
\[V_{I1}-V_{D1}-V_O=0\]
Odnosno\quad $0-V_D-V_O=0$. Kako je $V_o=RI_R$ onda je $I_R=-\dfrac{V_D}{R}=I_D$ pa bi struja kroz diodu bila negativna što je u suprotnosti sa početnom pretpostavkom da dioda vodi. Znači diode ne vode. Izlazni napon je $V_O=0$. Situacija se sigurno neće promeniti dok god su ulazni naponi manji od $V_{\gamma D}$.
\end{columns}
\end{frame}
